\section{Problem Setup and Damage Channels}
\label{sec:setup}

\subsection{Deletion Request and Retained-Behavior Scope}

Let $\theta_0$ denote the model before unlearning. A deletion request
specifies a forget set $\mathcal D_f$ and an enumerated universe
$\mathcal C$ of retained behaviors that should remain usable. An element
$x\in\mathcal C$ is a prompt--answer behavior on which retention can be
measured; it need not be a record from the original training set.

We split $\mathcal C$ at the semantic-group level into disjoint discovery
and audit folds,
$\mathcal C=\mathcal C_{\mathrm{disc}}\,\dot\cup\,\mathcal C_{\mathrm{audit}}$.
All profile-guided selection and protection use only
$\mathcal C_{\mathrm{disc}}$. Scores on $\mathcal C_{\mathrm{audit}}$ may be
computed before unlearning for prediction evaluation, but they are sealed
and cannot affect optimization, hyperparameter selection, checkpoint
selection, or stopping. Where a benchmark provides a native retain
partition, it is included in the sealed audit fold and denoted
$\mathcal A_{\mathrm{native}}$.

For a behavior $x$ and parameters $\theta$, let
\[
    \ell(x;\theta)
    =
    -\frac{1}{|y_x|}
    \sum_{j=1}^{|y_x|}
    \log p_\theta
    \bigl(y_{x,j}\mid p_x,y_{x,<j}\bigr)
\]
be the mean answer-token negative log-likelihood, with prompt tokens
masked. The realized damage to $x$ at a checkpoint $\theta_t$ of an
unlearning procedure is
\begin{equation}
    d_t(x)=\ell(x;\theta_t)-\ell(x;\theta_0),
    \label{eq:realized-damage}
\end{equation}
with positive values indicating harmful retention loss.

\subsection{Damage Channels}
\label{sec:setup-channels}

To first order in the update,
\[
    d(x)\approx g_x^{\top}\Delta\theta,
    \qquad
    |d(x)|\lesssim\lVert g_x\rVert_2\lVert\Delta\theta\rVert_2,
    \qquad
    g_x=\nabla_\theta\ell(x;\theta_0).
\]
This first-order coupling between a candidate's local gradient and the
realized update is the \emph{update-conditioned gradient interference}
mechanism named in the abstract. Gradient magnitude is then a
direction-free \emph{exposure envelope}; it
does not imply damage to every high-score candidate. Whether the envelope
ranks realized damage is an empirical interaction tested across objectives.
The method instantiates the envelope on a declared parameter block $B$
covering the parents' trainable parameters (Section~\ref{sec:method}).
The update is nevertheless structured: its gradients are set by where the
unlearning loss directly reads the model. Current objectives place their
forget-suppression signal at one of two readouts.

\paragraph{Output channel (loss-gradient).}
If the forget-suppression loss directly supervises outputs -- answer-token
likelihoods or preferences, as in GA, GradDiff, NPO, SimNPO, IdkDPO, and
GRU -- then $\Delta\theta$ is assembled from output-loss gradients. Trying
to predict each candidate's signed exposure $g_x^{\top}\Delta\theta$
requires anticipating the realized update direction, which is brittle over
a multi-step trajectory because that direction rotates with the model and
optimizer state. We instead measure the direction-free exposure envelope
$\lVert g_x\rVert_2$: how strongly the candidate loss can respond to small
parameter motion, before knowing which output-channel trajectory will run.

\paragraph{Representation channel.}
If the forget-suppression loss directly supervises internal hidden states,
as in RMU, RepNoise, and Circuit Breakers, the update is chosen to displace
activation geometry in a region occupied by the forget set. The matched
hypothesis is therefore \emph{proximity in hidden-state space}: candidates
whose representations share that region are more exposed to its
displacement. This is complementary to, rather than implied by, an
output-loss gradient norm.

\paragraph{Declaration, not post-hoc labeling.}
The channel of an objective is read off its loss definition before any
experiment is run -- a static inspection of which tensors enter the loss,
requiring no training run: token likelihoods or preferences imply the
output channel; hidden states imply the representation channel. For
objectives that mix terms (RepNoise adds an auxiliary output-side ascent
term), the declaration rule is the \emph{forget-suppression term}: the
loss component that drives removal decides the channel. Likewise, GRU
modifies the update rule but its losses read outputs, so it is declared
output-channel. We preregister the assignment (output: GA, GradDiff,
NPO, SimNPO, IdkDPO, GRU; representation: RMU, RepNoise, Circuit Breakers)
and never revise it against results. Two honesty clauses apply. First, the
dichotomy reflects where \emph{existing} objectives read the model, not a
law of nature; an objective reading some third quantity would define a
third channel, and the framework extends to it. Second, the channels are
poles of a spectrum rather than a partition of downstream effects: an
output-channel objective also perturbs representations incidentally, while
a representation objective ultimately changes outputs. The declaration is
about the objective's direct supervision signal, not every tensor in its
computation graph. Cross-channel signal is therefore measured rather than
ruled out; the confirmatory claim concerns the interaction between declared
channel and probe family.

\subsection{Prospective Prediction and the Interaction Endpoint}
\label{sec:setup-prediction}

A pre-unlearning predictor assigns each candidate a scalar score
$r(x;\theta_0,\mathcal D_f)$ before any target trajectory is run.
Predictive validity is evaluated against damage generated by independently
configured unlearning procedures that receive no access to $r$: within each
objective, we report Spearman correlation between the sealed score and
realized damage, plus tail-focused summaries.

The channel claim is not that some probe achieves high correlation -- it is
that \emph{which} probe family predicts flips with the objective's channel.
The preregistered primary endpoint is therefore a difference-in-differences
interaction. For a gradient-family probe $r_g$, a representation-family
probe $r_h$, an output-channel objective $\mathcal U_{\mathrm{out}}$, and a
representation-channel objective $\mathcal U_{\mathrm{rep}}$, with
$\rho(r,\mathcal U)$ the Spearman correlation of probe $r$ on damage from
$\mathcal U$:
\begin{equation}
    \Delta
    =
    \bigl[\rho(r_g,\mathcal U_{\mathrm{out}})-\rho(r_h,\mathcal U_{\mathrm{out}})\bigr]
    -
    \bigl[\rho(r_g,\mathcal U_{\mathrm{rep}})-\rho(r_h,\mathcal U_{\mathrm{rep}})\bigr].
    \label{eq:interaction}
\end{equation}
$\Delta>0$ with a candidate-bootstrap confidence interval excluding zero
supports channel matching. The double difference cancels both nuisance
factors that contaminate single comparisons: a probe that is uniformly
better, and an objective that is uniformly more predictable.

\subsection{Channel-Matched Protection and the Crossed Design}
\label{sec:setup-protection}

A score profile allocates a limited protection budget: using only the
discovery fold, a selector constructs a protected pool
$\mathcal P_{K_p}(r)\subseteq\mathcal C_{\mathrm{disc}}$ of the $K_p$
highest-ranked eligible behaviors, plus a matched low-score reference pool
$\mathcal R_0(r)$ (construction details in Section~\ref{sec:method}).
Protection methods are compared only at a common forgetting criterion
$\Gamma_f(\theta)\leq\tau_f$; a run that does not reach the criterion is a
reach failure, not a comparison point. Conditional on reach, the
confirmatory outcomes on the sealed audit are mean damage and upper-tail
damage $\operatorname{CVaR}_{.95}(\{d(x)\})$ -- the tail because that is
where collateral damage concentrates and where aggregate metrics hide it.

The operational endpoint is again an interaction, realized as a
preregistered crossed design: each parent objective (one per channel) is
run with each selector (matched family, mismatched family, random, none),
holding the parent, budgets, and repair procedure fixed. Success means the
channel-matched selector attains the lowest audit CVaR per parent. This
criterion deliberately does not ask our pipeline to outperform any
particular unlearner; it asks whether channel information -- and only
channel information -- moves protection where it helps.
